\subsection{补算：把whistler色散关系写成一个二乘二问题}
取扰动$\vct b=(b_x,b_y,0)e^{i(kz-\omega t)}$。因为$\partial_z=ik$，
\begin{equation}
 \grad\times\vct b=(-ikb_y,ikb_x,0),\qquad
 \vel_e={1\over\mu_0en}(ikb_y,-ikb_x,0).
\end{equation}
线性化感应方程只保留$\vel_e\times B_0\hat z$，舍去二阶扰动$\vel_e\times\vct b$。令$C=B_0/(\mu_0en)$，则
\begin{equation}
 -i\omega
 \begin{pmatrix}b_x\\b_y\end{pmatrix}
 =Ck^2
 \begin{pmatrix}0&-1\\1&0\end{pmatrix}
 \begin{pmatrix}b_x\\b_y\end{pmatrix}.
\end{equation}
第一个分量给$-i\omega b_x=-Ck^2b_y$，第二个给$-i\omega b_y=Ck^2b_x$。相互代入后
$\omega^2=C^2k^4$，故$|\omega|=Ck^2$。两个本征向量满足$b_y=\pm i b_x$，这就是两个相反旋转方向的圆偏振；给定正频率和传播方向时要选相容的一支。

电子惯性把$\vct b$替换成$\vct b_e=(1+d_e^2k^2)\vct b$，所以
\begin{equation}
 |\omega|={Ck^2\over1+d_e^2k^2},\qquad
 {C\over d_e^2}={eB_0\over m_e}=|\Omega_{ce}|.
\end{equation}
这也检查了单位：$C$是面积除以时间，乘$k^2$才是频率。

\subsection{频率、相速度、群速度不能混作一个数}
对正频率分支，
\begin{equation}
 v_{\rm ph}={\omega\over k}={Ck\over1+d_e^2k^2},\qquad
 v_{\rm g}={\dd\omega\over\dd k}
 ={2Ck\over(1+d_e^2k^2)^2}.
\end{equation}
无电子惯性时$\omega\propto k^2$、群速度为相速度的两倍；保留惯性后高$k$频率趋于常数，而相速度和群速度都下降。不能从“频率趋于$\Omega_{ce}$”推出“传播速度趋于常数”。

固定离子只是相对于研究时间与尺度的近似。该色散式不包括离子低频分支，也不包括位移电流和完整电磁高频响应。用它估计无惯性Hall显式算法的$\Delta t\propto\Delta x^2$很有用，但不能把它当作任意传播角、任意波数下的完整双流体色散关系。
